%% kapa: exam=1 round=26 year=2015 period=1 form=A subject=economics subject_ko=경제원론 seq=21 number=61
%% source: Q-Net 감정평가사 기출문제 (cst00309) artlSeq=5138001 file=extracted/1008/5138001/제26회_감정평가사 1차_기출/제26회_감정평가사 1차_기출/1교시_감평_A형.pdf
%% conversion: kordoc 4.24.51 -> tools/gen1.py (auto); review=proofread-sonnet; flags=none
\begin{question}{61}
거시경제변수에 관한 설명으로 옳지 않은 것은?
\begin{choices}
  \item GDP는 유량(flow)변수이다.
  \item GDP 디플레이터는 실질GDP를 명목GDP로 나눈 것으로 그 경제의 물가수준을 나타낸다.
  \item 기준연도의 명목GDP와 실질GDP는 같다.
  \item 외국인의 한국 내 생산활동은 한국의 GDP 산출에 포함된다.
  \item 소비, 투자, 정부지출(구입), 순수출이 GDP를 구성하는 네 가지 항목이다.
\end{choices}
\end{question}
\answer{2}
